\clearpage
\fancyhead[L]{\small CONNOISSANCE DES TEMS.}
\hypertarget{native-542-lower}{}\label{native-542-lower}
\begin{center}\small Göttingische gelehrte Anzeigen. No. 199. pp. 1978..1987. December 14, 1811.\\[3pt]\rule{.72\linewidth}{.3pt}\end{center}

\emph{Connoissance des tems ou des mouvemens célestes, à l'usage des astronomes et des navigateurs pour l'an} 1812, \emph{publiée par le bureau des longitudes. Juillet} 181c. Paris, published by \textsc{Courcier}. 415 pages in large octavo.

The table of longitudes and latitudes of the principal places on Earth has been considerably improved in this year's volume. Over a hundred new entries have been added, chiefly from \textsc{Monteiro}'s ephemerides of Coimbra and \textsc{Oltmanns}'s investigations into the geography of the New Continent, so that the total is now close to 1500. Mr. \textsc{Burckhardt} has also carefully reexamined the longitudes of 50 observatories or other \SourceLeaf{543}places in Europe of particular interest for the astronomical observations made there. We leave open whether the words il a comparé et calculé de nouveau \emph{toutes} les observations tant anciennes que modernes are to be taken quite literally. In any case, it would be desirable for Mr. \textsc{Burckhardt} to publish a fuller account of the details of this useful work. This time \textsc{Bouvard}'s observations at the Imperial Observatory in 1809 make up nearly half of the \emph{Additions}. The value of the observations of the Sun, Moon, planets and stars is well known. We were pleased to notice in this year's volume a number of observed transits of the Pole Star across the meridian. It is only a pity that they are all lower culminations, since this of course loses the many benefits that numerous observations of the Pole Star \emph{at both culminations taken together} can offer. There are quite a number of observations of Ceres and Pallas around the time of their opposition. There are likewise numerous eclipses of Jupiter's satellites, but only three incompletely observed stellar occultations, namely the emersion of $\gamma$ Scorpii on May 28 at $12^{\mathrm u}\,3^{\mathrm m}\,53^{\mathrm s}$ mean time, the emersion of $2\delta$ Tauri on Sept. 28 at $9^{\mathrm u}\,43^{\mathrm m}\,0^{\mathrm s}$, and the immersion of $1\delta$ Tauri on Oct. 25 at $18^{\mathrm u}\,17^{\mathrm m}\,15^{\mathrm s}2$. --- An easy way to calculate positions of the Moon approximately, by J. C. \textsc{Burckhardt}. Mr. B. uses the well-known Chaldean period of 18 years or 223 lunations, after which the arguments of the lunar inequalities return to nearly their former values. He calculates the changes in longitude and latitude, the most substantial of which are set out in five tables. From an ephemeris of the Moon's motion for a given year, one can thus calculate a similar ephemeris for the year one period later with very little effort and a fair degree of accuracy. We wish Mr. \textsc{Burckhardt} had also treated the change in horizontal parallaxes in a similar way. --- On a new way to improve pendulum clocks, by the same author. To prevent the pivots of the wheels from rusting, Mr. \textsc{Burckhardt} proposes fire-gilding them and then hardening them again. Putting this proposal into practice will best show whether anything is gained by it. --- Tables for the aberration, nutation and precession of \textsc{Maskelyne}'s 36 fundamental stars, by the same author. Solar nutation is combined with aberration. --- On the depression of mercury in barometer tubes due to capillarity, by \textsc{Laplace}. According to \textsc{Laplace}'s theory, the convex surface of the mercury, and hence its depression in the tube, is caused by the capillary action of an extremely fine film of moisture that remains on the glass surface even when the mercury is carefully boiled, rather than by the capillary action of the glass. The relation between the bore of the tube and the depression of the mercury depends on the shape of the surface. \textsc{Laplace} determined that shape by a laborious integration, based on mechanical quadrature, of the fundamental equation established in the supplement to Mécanique céleste (in our opinion, carrying the series expansion a little further might have made this somewhat easier). He thus calculated a table for practical use, which is given here. We wish that the height of the convex part of the mercury resulting from this calculation had also been included. --- On the solution given by \textsc{Gauss} (in the September 1808 issue of the Monatl. Correspondenz) to a problem in spherical astronomy, by \textsc{Delambre}. The problem is to determine at once the elevation of the pole at the place of observation, the clock offset and the instrument error from three equal altitudes of known stars. Mr. \textsc{Delambre}'s treatment of the problem here is \emph{essentially} the same as \textsc{Gauss}'s, but in our opinion it has lost in simplicity and elegance. Mr. \textsc{Delambre} has worked through the whole example of an observation with which \textsc{Gauss} illustrated the calculation at the time, going into great detail again. To \textsc{Gauss}'s result for the effect of observational errors on the accuracy of the latitude determination \SourceLeaf{544} \[ \mathrm{d}\varphi = 3{,}8077\Delta - 0{,}2884\Delta' + 3{,}5193\Delta'' \] he adds the remark Cette dernière formule prouve que dans la pratique la méthode n'aurait qu'une exactitude assez bornée. The only explanation for this misplaced judgment is that he has, rather incomprehensibly, overlooked that $\Delta$, $\Delta'$, $\Delta''$ are expressed in seconds of time and $\mathrm{d}\varphi$ in seconds of arc. Similarly, his further remark, dans une nuit, où l'on pourroit observer trois étoiles à la même hauteur, on verroit très probablement passer au méridien quelque étoile connue qui donneroit la latitude avec moins de peine et plus de précision, et l'heure de la pendule par une simple hauteur, shows that Mr. \textsc{Delambre} has again forgotten that this method is intended for all cases in which one cannot rely on the instrument for absolute altitudes. As for \emph{accuracy}, the remark is entirely unfounded. --- There follows some discussion of the familiar problem, almost completely useless in practice, of determining a single star's declinations and hour angle and the elevation of the pole all at once from three altitudes of that star. Mr. \textsc{Delambre} then returns once more to the short paper by \textsc{Gauss} that proposed observing two altitudes of two known stars to determine time and latitude, and that he discussed in the preceding volume of the Connaissance des tems. In what he says about this paper and the superiority of synthesis over analysis, we find further proof of the judgment we expressed in reviewing that volume of the Connoissance des tems, that Mr. \textsc{Delambre} had entirely misunderstood the purpose of the paper. The purpose, stated quite clearly in the paper itself, was to recommend a \emph{combination} of observations for determining the elevation of the pole. It is very useful for this under certain circumstances and, as far as the author knew then and knows to this day, had never been proposed for this purpose in such generality. That is why he called it a \emph{new} method. The \emph{calculation} from such observations, as \textsc{Gauss} also showed at the time, rests on a problem whose geometric solution has been known since \textsc{Tycho}'s day. After fully outlining that solution in a few words, the author went on to spend two pages showing that the same solution could also be found briefly and elegantly by purely analytical means. He did so only because he thought that those who enjoy analysis might welcome a derivation that is far from obvious and takes some skill. He never dreamed that anyone could take this to mean that all geometric treatment, whose value is well known and firmly established, was to be displaced. Mr. \textsc{Delambre}'s remarks about the relative merits of the analytical and geometric procedures are therefore quite out of place, even though they are for the most part fully justified. --- On the various methods astronomers have used to observe solar eclipses, by \textsc{Delambre}. A remarkable passage is quoted from \textsc{Appian}'s Astronomicum Caesareum, printed in 1540, in which colored glass is first proposed for observations of the Sun. Mr. \textsc{Delambre} has good grounds, however, for thinking it likely that \textsc{Appian} himself never tried to put his proposal into practice. --- Next comes a lengthy, 50-page excerpt from Gauss's \emph{Theoria motus corporum coelestium}. Mr. \textsc{Delambre} has excerpted most of its formulas, usually without the proofs, worked through the illustrative numerical examples again, though with less precision than in the work itself, and had the result printed here with occasional remarks of his own. This way of studying such works is much to be recommended, but such excerpts hardly seem suited to publication. For lack of space we can touch on only a few of the remarks here. Mr. \textsc{Delambre} applied to \emph{one} example the method of referring the calculation of the geocentric positions of the planets directly to the equator by means of rectangular coordinates, and he could not quite see the advantage of that \SourceLeaf{545}procedure over the usual one. Yet this goes entirely against the spirit of the method, which is intended solely for cases where \emph{many} geocentric positions are to be calculated. If Mr. \textsc{Delambre} had calculated a dozen planetary positions by this method instead of one, he would have found what so many other calculators have already found, that it takes less than half as much time and effort as the usual procedure. Mr. \textsc{Delambre} has also found for himself the four formulas of spherical trigonometry used so often in the Theoria motus, but he has failed to recognize their advantage over the usual procedure. Suppose, for example, that all the remaining parts of a spherical triangle are to be determined from two sides and the included angle. With the new formulas one has to look up a total of 12 logarithms at six different places in the sine tables. This also gives a check on the calculation, and the results are always precise and never ambiguous. With the usual procedure, which Mr. \textsc{Delambre} prefers, one must look up a total of 13 logarithms at 11 different places in the table if one wants a check on the calculation as well. One then obtains the third side from its sine, so that if the side is near a right angle, it is determined imprecisely and perhaps even ambiguously, although it is unambiguous in the problem itself. The advantage of the new formulas is therefore quite decisive and considerable enough when many such operations have to be carried out. It is surprising, then, that Mr. \textsc{Delambre} could have overlooked it. The wording Mr. \textsc{Delambre} uses on p.\,357 about \textsc{Gauss}'s second solution of the problem of determining the elements from two distances of a planet from the Sun, the included angle and the intervening time needs correction. He says that it requires \emph{very long} derivations that the author has not given and that it rests in part on theories of his own that he has not yet made public. This would lead one to believe that the solution to one of the most important problems in the whole work is presented in such a way that readers without considerable practice cannot manage it at all, and even experts cannot fully understand it. Both impressions are mistaken. The reviewer knows from many examples that even beginners have nowhere found any difficulty with the small, insignificant derivations that must always be left to the reader if the book is not to become unreadably tedious. The only connection between the solution and the author's own theories is through the computational shortcuts he himself used to construct the auxiliary table, which are by no means essential here. He had to reserve the development of those theories for another occasion, since this was no place for it. For the main equation in the great problem of determining the orbit from three geocentric positions, \textsc{Gauss} recommended the indirect solution as particularly convenient, but gave no further instructions on how to carry it out, partly because the techniques to be used are familiar enough in themselves and partly because the solution of \textsc{Kepler}'s problem given in the Theoria can serve to some extent as a model. Mr. \textsc{Delambre} nevertheless takes more than three times as long as necessary to solve the main equation numerically by the indirect method in the example he worked through again in the greatest detail. Though he himself still thinks his approach convenient enough, one can get further in two trials than Mr. \textsc{Delambre} does in seven. This shows that it might be worth the author's while to draw attention again, on a suitable occasion and in an appropriate place, to the little techniques to be used in the indirect solution. --- After the long review of the Theoria, in which Mr. \textsc{Delambre} confined himself to the first book and the first section of the second, there are also new remarks on the calculation of parallaxes and on the formulas of Messrs. \textsc{Olbers} and \textsc{Littrow}, by \textsc{Delambre}. He declares once again that he always regards solutions of problems based on \SourceLeaf{546}referring the positions of points in space to three rectangular coordinates and carrying out the derivations purely analytically as less simple than solutions that use spherical trigonometry alone. This is a very one-sided judgment with which few mathematicians will agree. \textsc{Olbers} had attributed to Mr. \textsc{Rohde} a formula expressing the parallax in longitude as a series in the sines of multiples of the Moon's distance from the nonagesimal. Mr. \textsc{Delambre} claims it here as his own invention, saying that he had already published it in the Connoissance des tems for 1793. It really belongs, however, to \textsc{Lagrange}, who had already published it in a somewhat more general form in the Nouveaux Mémoires de l'académie de Berlin, 1776, p. 231. --- A way of making a clock show sidereal time and mean time, by \textsc{Burckhardt}, is based on the approximate ratio of the two times, 51.79 to 49.82, which is off by only 4 seconds in a whole year. By the \emph{same author}, on the second comet of 1737, from observations in China that were published in the Monatl. Corresp. --- The meteorological observations at the Imperial Observatory in Paris in 1808 and the list of members of the Bureau of Longitudes close the volume, as usual.
